\section{A finite-coordinate toolkit}\label{stab:toolkit}

The estimates in this section apply to either coordinate model already
constructed. We write $Z$ for the chosen model: the finite-height
forest space $X$ of Proposition~\ref{presentation:identity}, or the
infinite-height coordinate predual $B_\infty$ of
Section~\ref{binf:definition}. In each case $J$ is the corresponding
primal segment completion and $\rho$ is its norm. The finite atomic
hulls and norming head tests are supplied by
Lemma~\ref{lem:finite-hull} in the forest and
Lemma~\ref{binf:finite-hull} in the infinite-height model.

For a finite ancestral set $H$, let $P_H$ be its coordinate projection
and put $Q_H=I-P_H$. These projections are contractive in either
model. Write $G_H$ for the minimal vertices of the complement of $H$.
Thus $G_H$ includes the roots of forest components missing $H$; if
$H$ is empty, it consists of all component roots. Its vertices are
incomparable, and their descendant subtrees partition the complement.
A crossing segment meets both $H$ and its complement and has a unique
first exterior vertex, its gate. A root of a component missing $H$
carries no crossing coefficient.

We use $d_S=\sum_{v\in S}e_v^*$ for a segment functional. The
coefficient estimate holds in both models:
\begin{equation}\label{stab:sign:coefficients}
 \left\|\sum_i a_i d_{S_i}\right\|
 \le\left(\sum_i a_i^2\right)^{1/2}
 \qquad(S_i\text{ pairwise disjoint}).
\end{equation}
Indeed, Cauchy--Schwarz bounds evaluation on $v\in J$ by the
coefficient norm times $\rho(v)$. An atom includes its
chosen representation
\[
 A=\sum_i c_i d_{S_i},\qquad S_i\text{ pairwise disjoint},
 \qquad \sum_i c_i^2\le1.
\]
A functional may have several such representations; all subsequent
statements concern the chosen one.

We will use the following common form of the cone estimates from the
forest construction and Lemma~\ref{binf:cones}. For a finite antichain
$G$ and vectors $u_g\in Z$ supported in its descendant subtrees,
\begin{equation}\label{stab:cone-bounds}
 \left(\sum_{g\in G}|u_g(e_g)|^2\right)^{1/2}
 \le\left\|\sum_{g\in G}u_g\right\|
 \le\left(\sum_{g\in G}\|u_g\|^2\right)^{1/2}.
\end{equation}
Tests supported on $G$ have $\rho$-norm equal to their Euclidean
norm, which proves the lower bound and the contraction of restriction
to an antichain. For the upper bound, combine finite disjoint segment
tests from the separate subtrees to obtain
$\sum_g\rho(P_{T_g}v)^2\le\rho(v)^2$, and apply
Cauchy--Schwarz. Choosing almost maximizing families suffices.
For countable $G$, finite partial sums give the same estimates whenever
$\sum_g\|u_g\|^2<\infty$; the upper bound makes those sums Cauchy.
If an existing vector is decomposed into its cone restrictions, its
finite cone sums converge by coordinate approximation, so the finite
estimates pass to that decomposition as well.

\subsection{The exterior term and two boundary comparisons}
\label{stab:coef:boundary}

Fix $p=P_Hp\in J$ with $\rho(p)\le1$. For a represented atom put
\[
 \begin{gathered}
 s_i=p(S_i),\qquad L(A)=1-A(p),\qquad D(A)=1-\sum_i s_i^2,\\
 E(A)=\sum_i(c_i-s_i)^2,\qquad A^\circ=\sum_i s_i d_{S_i}.
 \end{gathered}
\]
The evaluation deficit $L(A)$, squared-test slack $D(A)$, and
coefficient error $E(A)$ are nonnegative; also $A^\circ$ is an atom.
For each gate $g$, let $c_g(A),s_g(A)$ be the coefficient and head sum
of the crossing segment through $g\in G_H$, or zero if there is none.
There is at most one such segment in a disjoint family. Write
$c(A)=(c_g(A))_g$ and $s(A)=(s_g(A))_g$, and define
\[
 O(A)=\sum_{S_i\cap H=\varnothing}c_i d_{S_i},\qquad
 C(A)=\sum_{S_i\text{ crossing}}c_i d_{S_i\setminus H}.
\]
Thus $Q_HA=O(A)+C(A)$. The letters $O$ and $C$ denote the wholly
exterior part and the crossing part, respectively.

\begin{lemma}\label{stab:coef:pointwise}
For every represented atom, the exact identity is
\begin{equation}\label{stab:deficit-identity}
 2L(A)=1-\sum_i c_i^2+D(A)+E(A).
\end{equation}
In particular,
\begin{equation}\label{stab:coef:budget}
 D(A)+\sum_i(c_i-s_i)^2
 =1+\sum_i c_i^2-2A(p)\le2L(A).
\end{equation}
The full coefficient replacement satisfies
\begin{equation}\label{stab:toolkit:replacement}
 \max\{\|A-A^\circ\|,\|Q_H(A-A^\circ)\|\}
 \le\sqrt{E(A)}\le\sqrt{2L(A)}.
\end{equation}
Consequently,
\begin{equation}\label{stab:coef:exterior}
 \|O(A)\|\le\sqrt{2L(A)},\qquad
 \|c(A)-s(A)\|_2\le\sqrt{2L(A)}.
\end{equation}
At the gates one has
\begin{equation}\label{stab:coef:gate-coordinates}
 A(e_g)=c_g(A)+O(A)(e_g),\qquad
 \|(O(A)(e_g))_g\|_2\le\|O(A)\|.
\end{equation}
\end{lemma}

\begin{proof}
The identity is an expansion of the square; its right-hand inequality
uses $\sum_i c_i^2\le1$. Applying
\eqref{stab:sign:coefficients} to the coefficients $c_i-s_i$ and
using contractivity of $Q_H$ proves the replacement bound.
For the exterior segments $s_i=0$, so the
coefficient norm bound for disjoint segment functionals gives the first
inequality in \eqref{stab:coef:exterior}. The crossing indices form a
subset of the full coefficient array, giving the second. A crossing
segment contains its gate, and all other contributions there come from
$O(A)$, proving the coordinate identity. On any finite array supported
on the gate antichain, $\rho$ equals the Euclidean norm: a segment
meets that antichain at most once, and singleton tests give the reverse
inequality. Dual testing proves the final contraction, also for an
infinite gate set by taking finite subsets.
\end{proof}

\begin{lemma}\label{stab:coef:head-conflict}
Two represented atoms $A,B$, tested against the same $p$, satisfy
\begin{equation}\label{stab:coef:head-conflict-eq}
 2\sum_{g:s_g(A)s_g(B)<0}|s_g(A)s_g(B)|\le D(A)+D(B).
\end{equation}
There are two direct proofs: one discards a suffix of the longer head,
and the other splits the longer segment into two pieces.
\end{lemma}

\begin{proof}[Proof by truncation]
At an opposite-sign gate, both head segments end at the parent of that
gate. Their starting vertices lie on the same ancestor chain, so the
heads are nested. They are unequal because equal heads have equal
sums on $p$. Suppose $A$ has the longer head, with sum $s$, and the
shorter head of $B$ has sum $s'$. Replace the segment of $A$ by its
prefix ending just before the shorter head starts. This prefix is
nonempty, and its sum is $s-s'$. Its squared sum increases by
\[
 (s-s')^2-s^2=(s')^2-2ss'\ge2|ss'|.
\]
Perform these truncations at every opposite-sign gate where $A$ has
the longer head. Each original segment is used at most once, since it
has only one gate. The new segments are subsets of the old ones, so
they remain disjoint. Their total squared test sum is at most one,
and hence the total increase is at most $D(A)$. The other nesting
direction has budget $D(B)$. Adding proves the assertion.

This comparison can also be indexed first by the last interior vertex:
at most one crossing segment of an atom uses any given such vertex,
and the same nesting proof applies. Before using a cancellation of
coordinates, however, one must reindex by the \emph{first exterior}
vertex. Different children of one interior vertex are different
coordinates. Applying the same comparison after this reindexing gives
\eqref{stab:coef:head-conflict-eq} in precisely the form needed below.
\end{proof}

\begin{proof}[Proof by splitting]
Apply Lemma~\ref{lem:opposite-crossing} to the two underlying
segment families. In its notation,
$L_{\mathcal S}=D(A)$, $L_{\mathcal S'}=D(B)$, and the crossing
sums are $s_g(A),s_g(B)$. Its simultaneous two-piece cuts retain the
longer segment's original endpoint and give exactly
\eqref{stab:coef:head-conflict-eq}. The proof uses only nesting of
finite paths and the bound on their squared $p$-sums, so it applies
unchanged in either coordinate model. Gates that are untouched
component roots have zero crossing sums and contribute nothing.
\end{proof}

\begin{lemma}\label{stab:coef:original-products}
For the original crossing coefficients,
\begin{equation}\label{stab:coef:original-products-eq}
 \sum_g(-c_g(A)c_g(B))_+
 \le L(A)+L(B)+\sqrt{2L(A)}+\sqrt{2L(B)},
\end{equation}
where $r_+=\max(r,0)$.
\end{lemma}

\begin{proof}
The map $r\mapsto(-r)_+$ is one-Lipschitz. The difference between
the product sums formed with $c$ and with $s$ is therefore bounded by
\[
 \sum_g|c_g(A)-s_g(A)|\,|c_g(B)|
 +\sum_g|s_g(A)|\,|c_g(B)-s_g(B)|.
\]
The crossing arrays $c(B)$ and $s(A)$ have Euclidean norms at most
one. Cauchy--Schwarz and Lemma~\ref{stab:coef:pointwise} bound this
error by $\sqrt{2L(A)}+\sqrt{2L(B)}$.
Lemma~\ref{stab:coef:head-conflict} bounds the remaining sum by
$(D(A)+D(B))/2\le L(A)+L(B)$.
\end{proof}

\subsection{From scalar means to a tail norm}
\label{stab:coef:means}

Only finite probability spaces occur in the next two lemmas. This
allows every expectation, product measure, and maximizing segment to
be chosen before any passage to the completion.

\begin{lemma}\label{stab:coef:gate-bound}
Let $A$ have a finite law of represented atoms and suppose
$\mathbb E L(A)\le k$, where $k\ge0$. Write
\[
 \beta=\|\mathbb E c(A)\|_2,\qquad
 M_c=(\mathbb E|c_g(A)|)_g,\qquad
 M_s=(\mathbb E|s_g(A)|)_g.
\]
Then
\begin{align}
 \|\mathbb E s(A)\|_2&\le\beta+\sqrt{2k},
       \label{stab:coef:mean-head}\\
 \|M_s\|_2^2&\le(\beta+\sqrt{2k})^2+4k,
       \label{stab:coef:absolute-head}\\
 \|M_c\|_2&\le\sqrt{(\beta+\sqrt{2k})^2+4k}+\sqrt{2k},
       \label{stab:coef:via-head}\\
 \|M_c\|_2^2&\le\beta^2+4k+4\sqrt{2k}.
       \label{stab:coef:via-products}
\end{align}
The last inequality is obtained directly from the original products,
whereas the preceding one first controls the signs of the head sums.
\end{lemma}

\begin{proof}
The mean of $\|c-s\|_2$ is at most $\sqrt{2k}$, which proves
\eqref{stab:coef:mean-head}. Put
$a_g=\mathbb E(s_g)_+$ and $b_g=\mathbb E(-s_g)_+$.
For independent copies $A,B$, the expected opposite-sign product
sum is $2\sum_g a_gb_g$. Thus
Lemma~\ref{stab:coef:head-conflict} and
$\mathbb E D\le2k$ give
\[
 4\sum_g a_gb_g\le2\mathbb E D\le4k.
\]
Since $M_s=a+b$ and $\mathbb E s=a-b$, the same expansion gives
\begin{equation}\label{stab:signed-unsigned}
 \|M_s\|_2^2-\|\mathbb E s\|_2^2
 =2\mathbb E\sum_{g:s_g(A)s_g(B)<0}|s_g(A)s_g(B)|
 \le2\mathbb E D(A)\le4k.
\end{equation}
Together with \eqref{stab:coef:mean-head}, this proves
\eqref{stab:coef:absolute-head}. The coordinate inequality
$||c_g|-|s_g||\le|c_g-s_g|$ gives
$\|M_c-M_s\|_2\le\sqrt{2k}$, proving
\eqref{stab:coef:via-head}.

For the other route put $u_g=\mathbb E(c_g)_+$ and
$w_g=\mathbb E(-c_g)_+$. Lemma~\ref{stab:coef:original-products}
and Cauchy--Schwarz in the probability space give
\[
 2\sum_g u_gw_g
 =\mathbb E\sum_g(-c_g(A)c_g(B))_+
 \le2k+2\sqrt{2k}.
\]
Now $M_c=u+w$ and $\mathbb E c=u-w$, so
$\|M_c\|_2^2=\beta^2+4\sum_g u_gw_g$ gives
\eqref{stab:coef:via-products}.
\end{proof}

\begin{lemma}[Rooted mixtures and tail assembly]
\label{stab:coef:tail-assembly}\label{stab:sign:rooted}
Let $G$ be a finite antichain and, for each $g\in G$, let
\[
 f_g=\sum_j a_{g,j}d_{[g,w_{g,j}]},\qquad
 M_g=\sum_j|a_{g,j}|,
\]
where every sum is finite. Then
\begin{equation}\label{stab:sign:absolute-masses}
 \left\|\sum_g f_g\right\|\le\|(M_g)_g\|_2.
\end{equation}
If the nonzero coefficients have one common sign for each fixed $g$,
then
\begin{equation}\label{stab:sign:rooted-equality}
 \left\|\sum_g f_g\right\|
 =\left(\sum_g|f_g(e_g)|^2\right)^{1/2}.
\end{equation}
Consequently, adding another rooted mixture with the same sign at
each $g$ cannot decrease the norm.

In particular, for a finite law of represented atoms, put
$M_c=(\mathbb E|c_g(A)|)_g$. For any real random variable $\theta$
with $|\theta|\le1$,
\begin{equation}\label{stab:coef:tail-assembly-eq}
 \|\mathbb E(\theta C(A))\|\le\|M_c\|_2.
\end{equation}
If the law is an equal mixture of two laws with expectations
$\mathbb E_+,\mathbb E_-$, then
\begin{equation}\label{stab:coef:conditional-assembly}
 \|\mathbb E_\pm C(A)\|
 \le\|(\mathbb E_\pm|c_g(A)|)_g\|_2\le2\|M_c\|_2.
\end{equation}
The same estimates pass to norm limits of finite rooted mixtures
when their absolute-mass vectors converge in $\ell_2$. In particular,
countably many roots are permitted when their masses are square
summable and the series at each root is absolutely summable.
\end{lemma}

\begin{proof}[Proof by convex products]
Ignore roots with $M_g=0$. At each remaining root, select $j$ with
probability $|a_{g,j}|/M_g$ and choose the signed segment
$\operatorname{sgn}(a_{g,j})d_{[g,w_{g,j}]}$. Make these choices
independently at the finitely many roots. Put $r=\|(M_g)_g\|_2$.
If $r=0$, all mixtures vanish. Otherwise multiply the selected
segment at $g$ by $M_g/r$. The selected segments lie in incomparable
subtrees and are disjoint, and their squared coefficients sum to one.
Every resulting functional is therefore an atom. Its expectation is
$(\sum_g f_g)/r$, which proves \eqref{stab:sign:absolute-masses}.

For the crossing expectation in \eqref{stab:coef:tail-assembly-eq},
the contribution at $g$ divided by $M_c(g)>0$ is a convex combination
of signed segments starting at $g$, with zero allowed to account for
$|\theta|<1$. The same finite product construction applies; choosing
zero simply omits that segment. Apply it separately to either
conditional law. Since
$\mathbb E_\pm|c_g|\le2\mathbb E|c_g|$ coordinatewise, it also
proves \eqref{stab:coef:conditional-assembly}.

Under the one-sign assumption, $M_g=|f_g(e_g)|$.
Antichain testing gives the reverse inequality in
\eqref{stab:sign:rooted-equality}. Adding same-sign coefficients
increases every absolute root coordinate, proving the asserted
monotonicity. This conclusion concerns these rooted mixtures; it does
not assert lattice monotonicity of the norm on $Z$.

For the stated extensions, apply the finite inequalities first.
Their upper bounds show that truncating absolutely summable
coefficients at each root and then square-summable root masses gives
a Cauchy sequence in $Z$. Root coordinates converge by continuity,
so the equalities pass to the limit as well. No infinite product is
needed to produce an atom.
\end{proof}

\begin{proof}[Proof by disjoint test segments]
For $v\in c_{00}$, let $m_g(v)$ be the largest $|d_R(v)|$ among
the rooted segments appearing at $g$, or zero if there are none.
Choose a maximizing segment for every nonzero $m_g(v)$. They lie
in incomparable subtrees and are disjoint, so
$\sum_gm_g(v)^2\le\rho(v)^2$. Consequently,
\[
 \left|\sum_g f_g(v)\right|
 \le\sum_g M_gm_g(v)\le\|(M_g)_g\|_2\rho(v).
\]
Density extends this bound to $J$. For a finite atom law use the
crossing tails in its support in the same calculation:
\[
 |\mathbb E(\theta C(A))(v)|
 \le\sum_g\mathbb E|c_g|\,m_g(v)
 \le\|M_c\|_2\rho(v).
\]
The conditional argument is identical, and the root tests used in
the first proof again give the one-sign equality.

One may maximize over all rooted segments in a fixed finite ancestral
set, or take the supremum over all such segments in a subtree.
In the latter case, choose almost maximizing segments for finitely
many roots before taking the limit. This proves the same estimate
without assuming global attainment. The norm-limit statements then
follow by the finite truncations described above.
\end{proof}
